%% {{CODE}} article template -- {{PUBNAME}}
%%
%% Limits: no more than 5 pages in this template, including figures, tables and
%% references (about 4,500 words of main text; Expert Perspectives 3 pages);
%% an abstract of up to 150 words; no more than 6 figures and tables combined;
%% no more than 3 displayed equations; no more than 15 references; no
%% appendices. Manuscripts over these limits are returned before review.
%%
%% Submit this file compiled with the `anonymous` option, which removes the
%% author block for double-blind review. Give author names, affiliations,
%% biographies and acknowledgements in titlepage.tex instead, and submit that
%% separately. `lineno` adds line numbers for the reviewers; keep it for
%% submission.
\documentclass[magazine,anonymous,lineno,{{CITEOPT}}]{nextgeniq}

\publication{{{PUBNAME}}}{{{CODE}}}
\articletype{Technical Feature} % or Review Article, Tutorial, Case Study, Expert Perspective
\shorttitle{A short running title of up to 50 characters}

\title{Title of the Article: Short, Clear and Engaging}

% Author details are hidden by the `anonymous` option. Fill them in anyway:
% the same block is used for the accepted version.
\author[1]{First Author}
\author[2]{Second Author}
\author[1]{Third Author}
\affil[1]{Department, University, City, Country}
\affil[2]{Department, Institution, City, Country}
\correspondence{First Author, \texttt{first.author@example.org}}

\keywords{keyword one; keyword two; keyword three; keyword four}
% Article history: completed by the editorial office. Leave these as they are.
\received{DD Month YYYY}
\revised{DD Month YYYY}
\accepted{DD Month YYYY}
\published{DD Month YYYY}

\begin{document}

\maketitle

\begin{abstract}
In no more than 150 words, say what the article is about, why it matters now,
and what readers will take away. Do not cite references or use equations in
the abstract.
\end{abstract}

\section{Introduction}

Open with the problem and why it matters now, before any technical detail.
Write for readers outside your specialism: researchers, engineers,
practitioners and decision-makers. Close the introduction with what the
article offers the reader. Cite sources by number \citep{smith2024}, in order
of first citation.

Define every abbreviation and acronym at first use in the text, even if it is
already defined in the abstract, and avoid abbreviations in the title. Use
italics for emphasis; do not underline. Do not change font sizes or spacing to
fit the page limit.

\section{Background}

Give the context a non-specialist reader needs: the state of practice, earlier
approaches and the gap this article addresses \citep{lee2023,garcia2022}.

\begin{sidebar}{Sidebar: Key Terms}
Put background that specialists can skip in a sidebar of no more than 300
words, and refer to it in the text. A sidebar may hold definitions, a short
history, or a summary of a standard.
\end{sidebar}

\section{Approach}

Explain each idea in words first, using a figure where it helps
(Figure~\ref{fig:example}). Describe the system, method or framework at the
level a practitioner needs to understand and apply it.

\subsection{Equations}

Keep equations to a minimum: no more than three displayed equations, and only
where they are essential. Number them consecutively, as in
Equation~\eqref{eq:example}. Define every symbol before or immediately after
the equation.
\begin{equation}
  \hat{y} = \arg\max_{y} \; P(y \mid x)
  \label{eq:example}
\end{equation}

\subsection{Algorithms}

Number algorithms and give each a short title, as in
Algorithm~\ref{alg:example}. Set them off with rules above the title and
below the last line, and prefer clear pseudocode to program code.

\begin{algorithmbox}{A short, descriptive title}\label{alg:example}
  \item \textbf{Input:} the data set $D$ and the threshold $\tau$
  \item \textbf{Output:} the selected items $S$
  \item $S \gets \emptyset$
  \item \textbf{for each} item $x \in D$ \textbf{do}
  \item \quad \textbf{if} $\mathrm{score}(x) \ge \tau$ \textbf{then} add $x$ to $S$
  \item \textbf{end for}
  \item \textbf{return} $S$
\end{algorithmbox}

\begin{figure}[t]
  \centering
  \includegraphics[width=\linewidth]{example-image}
  \caption{Figure captions go below the figure. Use figures to explain ideas,
  such as architectures, workflows and comparisons. No more than six figures
  and tables combined. Use \texttt{figure*} for a figure across both
  columns.}
  \label{fig:example}
\end{figure}

\section{Results}

Present only the key results that support the article's message, usually in
one or two figures or a single table such as Table~\ref{tab:results}.
Parameter sweeps, ablations and long result tables belong in a journal
article.

\begin{table}[t]
  \centering
  \caption{Table captions go above the table. Use horizontal rules only.}
  \label{tab:results}
  \begin{tabular}{lccc}
    \toprule
    Method   & Precision & Recall & F1 \\
    \midrule
    Baseline & 0.81      & 0.77   & 0.79 \\
    Proposed & 0.88      & 0.85   & 0.86 \\
    \bottomrule
  \end{tabular}
\end{table}

\subsection{Preparing figures and tables}

\begin{itemize}
  \item Supply photographs and colour or greyscale figures at 300\,dpi or
  more, and line art at 600\,dpi, or use vector graphics (PDF, EPS). Submit
  each figure as a separate file as well.
  \item Size figures to one column (about 85\,mm) or the full page width
  (about 175\,mm). Text in figures should be 8--10\,pt at print size.
  \item Label axes with words and units in parentheses, such as
  ``Latency (ms)'', not symbols or units alone.
  \item Label the parts of a multipart figure (a), (b), (c) beneath each part.
  Do not draw borders around figures or put captions inside them.
  \item Cite every figure and table in the text, and place it near its first
  mention.
\end{itemize}

\section{Discussion}

Explain what the results mean in practice, state the limitations honestly,
and describe what remains unsolved and where research and practice should go
next.

\section{Conclusion}

Summarise the take-away message in a short paragraph. Do not repeat the
abstract; say why the work matters and what readers should do with it.

\section*{Acknowledgements}

\textit{Leave this section empty for review: acknowledgements can identify the
authors. Give them in titlepage.tex; they are added at acceptance.}

\begin{declarations}
  \declaration{Funding} Name the funders and grant numbers, or state that the
  work received no external funding. Anonymise funder details that would
  identify the authors during review.
  \declaration{Conflicts of interest} Declare any competing interests, or
  state that there are none.
  \declaration{Data availability} Where the article reports data, state where
  it can be found, or why it cannot be shared.
  \declaration{Ethics approval} Give the approving body and reference for
  work involving human participants, or state that approval was not
  required.
  \declaration{Use of generative AI} Describe any substantive use of
  generative AI tools, or state that none was used.
\end{declarations}

% No more than 15 references, preferably recent and peer-reviewed. Include a
% DOI wherever one exists.
\bibliography{references}

\end{document}
